\section{Consultas y Analítica Avanzada}

\subsection{Traducción del lenguaje natural a SQL para el histórico de movimientos}

Para que el usuario pueda preguntar por sus gastos pasados de forma natural, 
utilizamos \textit{Text-to-SQL} a través de la herramienta \texttt{consultar\_movimientos}. 
De esta forma, dejamos que el modelo traduzca la pregunta directamente a una consulta en SQLite.

Para que el modelo no se equivoque interpretando los datos, le pasamos unas reglas muy claras en el 
\textit{System Prompt}:

\begin{itemize}

    \item \textbf{Tratamiento de signos:} En nuestra base de datos, los gastos se guardan en negativo. 
    Para sumar lo gastado y que la cifra no confunda al usuario, le explicamos al modelo que tiene que usar 
    \texttt{SUM(-importe)} y mostrar siempre el resultado en positivo.
    
    \item \textbf{Gestión temporal:} Le indicamos cómo usar funciones nativas como \texttt{strftime} (ya que el manejo de fechas en SQLite puede generar errores si no se unifican los formatos).
    
    \item \textbf{Autocorrección:} Si el modelo genera un SQL inválido que produce un error en SQLite, el sistema 
    no se rompe. La capa de base de datos intercepta el error y se lo devuelve al agente en el siguiente turno. 
    De esta forma, el modelo se da cuenta de su propio fallo, corrige la sentencia y lo vuelve a intentar automáticamente 
    (hasta un máximo de dos veces para no quedarse en bucle).
    \begin{figure}[!ht]
\centering
\resizebox{\textwidth}{!}{%
\begin{tikzpicture}[
    font=\small,
    llm/.style={draw=ugrred, fill=ugrred!10, rounded corners=4pt, align=center, minimum width=3.8cm, minimum height=1.6cm, thick},
    backend/.style={draw=subsectionblue, fill=subsectionblue!10, rounded corners=4pt, align=center, minimum width=3.8cm, minimum height=1.6cm, thick},
    db/.style={draw=gray!70!black, fill=gray!10, rounded corners=4pt, align=center, minimum width=3.8cm, minimum height=1.6cm, thick},
    arrow/.style={-{Stealth[length=2.5mm]}, thick, draw=gray!80!black},
    lbl/.style={font=\scriptsize, inner sep=4pt, text=gray!80!black} % Sin fondo blanco para no cortar líneas
]

% 1. Nodos principales
\node[llm]     (llm)    at (0, 0) {\textbf{Modelo de Lenguaje}\\(Traduce a SQL)};
\node[backend] (api)    at (7.5, 0) {\textbf{Capa de Control}\\(\texttt{ejecutar\_sql\_seguro})};
\node[db]      (sqlite) at (15, 0) {\textbf{Base de Datos}\\(SQLite \texttt{mode=ro})};

% 2. Flechas directas (Ida y vuelta a la BD)
\draw[arrow] ([yshift=4mm]llm.east) -- node[lbl, above=1pt] {1. \texttt{SELECT ...}} ([yshift=4mm]api.west);
\draw[arrow] ([yshift=4mm]api.east) -- node[lbl, above=1pt] {2. Ejecución} ([yshift=4mm]sqlite.west);
\draw[arrow] ([yshift=-4mm]sqlite.west) -- node[lbl, below=1pt] {3. Resultados / Error} ([yshift=-4mm]api.east);

% 3. Ruta de Error (Bucle amplio por arriba, texto ENCIMA de la línea)
\draw[arrow, draw=ugrred, rounded corners=12pt] 
    ([xshift=-6mm]api.north) -- 
    (6.9, 2.6) -- 
    node[lbl, above=1pt, text=ugrred] {\textbf{Error sintáctico} $\rightarrow$ 4a. Feedback al modelo (Máx. 2 reintentos)} 
    (0.6, 2.6) -- 
    ([xshift=6mm]llm.north);

% 4. Ruta de Éxito (Bucle amplio por abajo, texto DEBAJO de la línea)
\draw[arrow, draw=green!60!black, rounded corners=12pt] 
    ([xshift=-6mm]api.south) -- 
    (6.9, -2.6) -- 
    node[lbl, below=1pt, text=green!60!black] {\textbf{Éxito (Filas)} $\rightarrow$ 4b. Datos truncados al límite (Máx. 50)} 
    (0.6, -2.6) -- 
    ([xshift=6mm]llm.south);

\end{tikzpicture}%
}
\caption{Flujo \textit{Text-to-SQL} y mecanismo de autocorrección.} %Si la consulta generada por el modelo contiene errores o palabras prohibidas, el \emph{backend} interrumpe la respuesta y devuelve el error al modelo para que genere una nueva instrucción corregida.}
\label{fig:text_to_sql}
\end{figure}
\end{itemize}

\newpage
Con este diseño, el modelo actúa como un traductor ágil entre el usuario y los datos, mientras que la ejecución 
segura y la gestión de fallos siguen controladas por el \emph{backend}.

\subsection{Enfoque híbrido: heurística y estadística}

Hay preguntas que un modelo de lenguaje pequeño no sabe deducir solo escribiendo SQL. Averiguar qué pagos son periódicos 
o predecir cuánto dinero nos va a sobrar a final de mes exige hacer cálculos sobre todo el historial y buscar patrones 
estadísticos.   

Por eso, apostamos por un enfoque híbrido: el análisis estadístico está programada de forma determinista en Python 
dentro del módulo \texttt{analitica.py}, y el modelo de lenguaje solo entra en juego al final, para entender qué ha pedido 
el usuario y redactarle el resultado.

\subsubsection{Detección de pagos recurrentes y suscripciones}

En la base de datos no hay ninguna columna que nos diga si un pago es una suscripción mensual
o un recibo fijo; es algo que hay que inferir analizando el histórico.

Para detectarlos, nuestro algoritmo no se fija en que el importe sea siempre idéntico 
porque recibos como la luz o el agua varían cada mes, sino en la regularidad de los 
días que pasan entre un cobro y otro.

El procedimiento es el siguiente: calculamos la mediana de los intervalos entre los pagos 
de un mismo comercio y medimos cuánto se desvían de ella. Si ese factor de irregularidad es menor o igual 
a 0,06 (valor elegido empíricamente), y la mediana encaja en una cadencia conocida (mensual, trimestral, anual, etc.) con un margen de 
tolerancia del 20\%, confirmamos que es un pago recurrente.

Esto nos permite identificar desde la cuota de Netflix hasta el seguro del coche que llega una vez al año. 
Además, al revisar todo el histórico, este mismo algoritmo detecta automáticamente si alguna de estas 
suscripciones ha subido de precio recientemente.

\subsubsection{Proyecciones de gasto y analítica predictiva}

Para estimar el gasto total de un mes no resulta suficiente con extrapolar
directamente el gasto acumulado hasta el momento. Esta aproximación puede
producir resultados poco representativos al comienzo del mes, especialmente
cuando existen pagos fijos que todavía no se han realizado.

Para que la previsión sea realmente útil, decidimos dividirla en dos bloques:

\begin{itemize}

    \item \textbf{Gastos fijos:} Como ya hemos detectado los pagos recurrentes en el paso anterior, 
    sabemos cuánto cuestan al mes. Los sumamos enteros a la proyección directamente, independientemente 
    de si ya se han cobrado este mes o están pendientes. Hacerles una regla de tres rompería la estimación.

    \item \textbf{Gasto variable:} Aquí sí hacemos una estimación, mezclando el ritmo de gasto de este mes con la 
media histórica de los últimos 12 meses. El peso de cada parte se ajusta dinámicamente según el día en el que estemos: 
el día 3 pesa más la historia, y el día 25 nos fiamos más del ritmo actual

\end{itemize}

Para afinar y validar este modelo usamos \textit{backtesting} sobre 23 meses de histórico real, lo que nos permitió 
medir exactamente cuánto nos equivocábamos cada día del mes según la categoría de gasto

Los resultados del análisis se resumen antes de enviarlos al modelo de
lenguaje. En lugar de proporcionarle únicamente los valores numéricos
calculados, el backend puede generar conclusiones como ``fiabilidad alta'' o
``tendencia por encima de la media''. De esta forma, el modelo recibe
información ya procesada y se encarga principalmente de presentarla al usuario.